% ===========================================================================
\chapter{Architecture de l'application}
\label{chap:archi-app}
% ===========================================================================

\begin{constat}{état de ce chapitre}
Ce chapitre documente les choix de conception arrêtés. L'implémentation
correspondante relève de la phase~2 et sera documentée au
chapitre~\ref{chap:dev-securise}. Les diagrammes définitifs seront produits en
phase~1 avec le modèle de menaces.
\end{constat}

\section{Choix du domaine fonctionnel}

Le cas d'étude retenu est une \textbf{place de marché d'articles vintage} :
des vendeurs publient des objets d'occasion uniques, des acheteurs les
acquièrent.

Ce domaine n'a pas été choisi pour son intérêt fonctionnel mais pour la
richesse de sa surface de sécurité. Le tableau~\ref{tab:surface} justifie
chaque fonctionnalité retenue par la catégorie de risque qu'elle permet de
démontrer.

\begin{table}[H]
\centering
\caption{Justification des fonctionnalités par la surface de sécurité créée}
\label{tab:surface}
\small
\begin{tabular}{@{}p{3.9cm}p{4.0cm}p{6.0cm}@{}}
\toprule
\textbf{Fonctionnalité} & \textbf{Surface} & \textbf{Risque démontré} \\
\midrule
Inscription, connexion & Authentification
& Authentification défaillante, bourrage d'identifiants, stockage faible des
mots de passe \\
\addlinespace
Rôles acheteur, vendeur, administrateur & Autorisation
& Élévation de privilèges, contrôle d'accès au niveau des fonctions \\
\addlinespace
Gestion par le vendeur de ses seules annonces & Propriété des ressources
& \textbf{IDOR / BOLA} --- la démonstration centrale du projet \\
\addlinespace
Téléversement d'images & Fichiers non fiables
& Téléversement malveillant, traversée de chemin, exécution côté serveur \\
\addlinespace
Recherche d'annonces & Entrée utilisateur en requête
& Injection SQL, énumération, moissonnage \\
\addlinespace
Commande, paiement simulé & Logique métier à valeur
& Abus de logique métier, \textbf{double vente d'un objet unique} \\
\addlinespace
Avis sur commande & Contenu stocké
& XSS stocké \\
\addlinespace
Modération & Interface privilégiée
& Contrôle d'accès sur fonction d'administration \\
\addlinespace
Génération de description par IA & Intégration d'un LLM
& Injection d'invite, traitement non sûr de la sortie, abus de consommation \\
\bottomrule
\end{tabular}
\end{table}

\begin{decision}{périmètre fonctionnel gelé}
Sont explicitement hors périmètre : le traitement réel des paiements, l'envoi de
courriels, l'authentification fédérée et la messagerie temps réel.

Aucun de ces éléments n'apporte une catégorie de risque que le périmètre retenu
ne couvre déjà, et chacun consommerait du temps appartenant à la pipeline.
Toute fonctionnalité ajoutée devra en remplacer une autre.
\end{decision}

\section{Composants et choix techniques}

\begin{table}[H]
\centering
\caption{Composants de l'application}
\label{tab:composants}
\small
\begin{tabular}{@{}llp{6.4cm}@{}}
\toprule
\textbf{Composant} & \textbf{Technologie} & \textbf{Justification} \\
\midrule
API métier & Express.js, TypeScript
& Écosystème imposé par le sujet (Projet 3) ; typage statique \\
\addlinespace
Base de données & PostgreSQL
& Transactions et verrouillage de ligne, nécessaires au contrôle de double
vente \\
\addlinespace
Accès aux données & Prisma
& Requêtes paramétrées par défaut, \asvsc{1.2.4} \\
\addlinespace
Validation & Zod
& Validation par schéma en couche de service de confiance, \asvsc{2.2.2} \\
\addlinespace
Mots de passe & Argon2id
& Fonction de dérivation approuvée, \asvsc{11.4.2} \\
\addlinespace
En-têtes HTTP & Helmet
& En-têtes de sécurité du navigateur, \asvsc{3.4.x} \\
\addlinespace
Front-end web & React, TypeScript, Vite
& Partage du contrat d'API et des types avec le back-end \\
\addlinespace
Client mobile & React Native (Expo)
& \emph{Reporté en phase 2 du projet} ; réutilise la même API \\
\addlinespace
Fournisseur LLM & Google Gemini (offre gratuite)
& Aucun coût ; quota contraint, ce qui rend le risque \llm{10} concret \\
\bottomrule
\end{tabular}
\end{table}

\captureaplacer{figures/architecture-composants.png}{Diagramme de composants de
l'application : front-end web, client mobile (phase 2), API REST, base de
données, stockage d'objets, et fournisseur LLM identifié comme frontière de
confiance externe. À produire en phase 1.}

\section{Interface de programmation REST}

Le tableau~\ref{tab:api} présente la surface d'API prévue. Deux points
d'entrée sont désignés dès la conception comme cibles des démonstrations de
vulnérabilité du chapitre~\ref{chap:vulnerabilites}.

\begin{table}[H]
\centering
\caption{Surface d'API REST}
\label{tab:api}
\small
\begin{tabular}{@{}llp{5.9cm}@{}}
\toprule
\textbf{Point d'entrée} & \textbf{Accès} & \textbf{Contrôle principal} \\
\midrule
\texttt{POST /api/auth/register}   & public & Politique de mot de passe, limitation de débit \\
\texttt{POST /api/auth/login}      & public & Message d'erreur générique, anti-automatisation \\
\texttt{POST /api/auth/refresh}    & cookie & Rotation et détection de réutilisation \\
\texttt{POST /api/auth/logout}     & authentifié & Invalidation côté serveur \\
\texttt{GET /api/listings}         & public & Requête paramétrée, pagination \\
\texttt{GET /api/listings/:id}     & public & Champs publics uniquement \\
\texttt{POST /api/listings}        & vendeur & Validation, garde d'affectation de masse \\
\texttt{PUT /api/listings/:id}     & \textbf{propriétaire} & \textbf{Cible de la démonstration V5} \\
\texttt{DELETE /api/listings/:id}  & propriétaire & Propriété revérifiée côté serveur \\
\texttt{POST /api/listings/:id/images} & propriétaire & Octets magiques, ré-encodage \\
\texttt{POST /api/orders}          & acheteur & Verrou transactionnel \\
\texttt{GET /api/orders}           & \textbf{propriétaire} & \textbf{Cible de démonstration BOLA} \\
\texttt{POST /api/reviews}         & acheteur & Assainissement en sortie \\
\texttt{GET /api/admin/listings}   & administrateur & Contrôle d'accès par fonction \\
\texttt{POST /api/ai/describe}     & vendeur & Gardes d'injection, plafond de jetons \\
\bottomrule
\end{tabular}
\end{table}

\section{Modèle de données}

Les entités prévues sont : \texttt{User}, \texttt{Listing},
\texttt{ListingImage}, \texttt{Order}, \texttt{Review} et
\texttt{RefreshToken}.

Une contrainte de conception mérite d'être soulignée : l'entité
\texttt{Listing} doit porter un état permettant de distinguer un article
disponible d'un article réservé ou vendu. Cet état est le support du contrôle de
non-duplication de vente (\asvsc{2.3.4}), lequel constitue la particularité
d'une place de marché d'objets uniques par rapport à un commerce de biens en
stock.

\section{Classification des données}

\begin{constat}{à compléter en phase 1}
La classification des données et les exigences de protection associées à chaque
niveau constituent une exigence documentaire de l'ASVS (\asvsc{14.1.1} et
\asvsc{14.1.2}). Elles seront produites en phase~1, en même temps que le modèle
de menaces, et reportées ici.
\end{constat}

\section{Architecture indifférente au client}

L'API est conçue dès l'origine pour servir indifféremment le front-end web et
le futur client mobile. Ce choix a une conséquence pratique sur le report de
l'application mobile exposé au chapitre~\ref{chap:introduction} : l'ensemble des
contrôles côté serveur, le modèle de menaces, la pipeline, la suite de tests
d'abus et les exigences ASVS côté serveur sont réutilisés sans modification.

Les travaux propres au mobile se limitent au client lui-même et aux risques qui
lui sont spécifiques --- stockage local, épinglage de certificat, intégrité de
l'artefact --- recensés au chapitre~\ref{chap:limites}.
